\chapter{Trichomoniasis}

\section*{What Is This Test?}
Trichomoniasis testing detects \textit{Trichomonas vaginalis}, a flagellated, motile, anaerobic protozoan parasite and the most common non-viral STI worldwide, with an estimated 3.7 million prevalent infections in the United States. \textit{T. vaginalis} infects the squamous epithelium of the vagina, urethra, and paraurethral glands in women, and the urethra, prostate, and epididymis in men. In women, trichomoniasis causes a diffuse, malodorous, yellow-green, frothy vaginal discharge with vaginal erythema, "strawberry cervix" (colpitis macularis/punctate hemorrhages on the cervix, visible in only 2--5\% of cases on naked-eye exam but frequently seen on colposcopy), dysuria, and pruritus. However, 70--85\% of infected individuals are asymptomatic. In men, it is usually asymptomatic (75\%) but may cause non-gonococcal urethritis (NGU), epididymitis, or prostatitis. Trichomoniasis is associated with adverse outcomes including increased risk of HIV acquisition and transmission (2- to 3-fold), preterm birth, premature rupture of membranes, low birth weight, and enhanced susceptibility to HPV infection and cervical neoplasia. \textit{T. vaginalis} is not nationally notifiable and is often underdiagnosed because reliance on wet-mount microscopy misses 40--60\% of cases.

\section*{Alternative Names}
Trich Test; T. vaginalis NAAT; Trichomonas PCR; Wet Mount; Trichomonas Culture; TV NAAT; OSOM Trichomonas Rapid Test; Affirm VPIII

\section*{Principle}
NAATs are the gold standard with the highest sensitivity. FDA-cleared platforms include: Hologic Aptima \textit{Trichomonas vaginalis} assay (targets 18S rRNA by transcription-mediated amplification, sensitivity 95--100\%, specificity >99\%), Cepheid GeneXpert TV (real-time PCR, sensitivity 96--100\%), and BD MAX CT/GC/TV (multiplex real-time PCR, sensitivity 95--98\%). NAATs are FDA-cleared for vaginal swabs (clinician-collected or self-collected), endocervical swabs, and female urine; the Aptima assay is also cleared for male urine. Results available in 1--4 hours. Wet-mount microscopy: A drop of vaginal discharge is mixed with a drop of normal saline on a glass slide, covered with a coverslip, and examined immediately under 400x magnification for motile, flagellated, pear-shaped trichomonads (jerky, tumbling motility). Sensitivity is only 40--60\% compared to NAAT, and motility is lost within 10--20 minutes of specimen collection. The wet mount is the most accessible test but should not be used to exclude trichomoniasis. OSOM Trichomonas Rapid Test (antigen detection): Immunochromatographic dipstick test detecting \textit{T. vaginalis} membrane proteins; sensitivity 82--90\%, specificity >97\%. Results in 10 minutes. Culture: Modified Diamond's medium or InPouch TV system; sensitivity 80--90\%; requires 2--5 days incubation. Now largely supplanted by NAAT. Papanicolaou (Pap) smear identification of \textit{T. vaginalis}: Incidental identification on a routine Pap smear has low sensitivity (~60\%) and poor specificity (many false positives); should NOT be used as a diagnostic test but positive findings warrant confirmatory NAAT.

\section*{History}
\textit{Trichomonas vaginalis} was first described by Alfred Donné in 1836 in vaginal secretions. For most of the 20th century, diagnosis relied on wet-mount microscopy, which remained the standard despite limited sensitivity. Metronidazole was introduced in the 1960s and remains first-line therapy. The InPouch TV culture system was developed in the 1980s. Antigen detection tests (OSOM) were approved in the 2000s and improved point-of-care sensitivity. NAATs for \textit{T. vaginalis} were introduced in the 2010s and became FDA-cleared; the CDC's 2021 STI Treatment Guidelines strongly recommend NAAT as the diagnostic method of choice. Despite these advances, trichomoniasis remains underdiagnosed because screening is not as widespread as for chlamydia and gonorrhea, and many laboratories do not employ NAAT for \textit{T. vaginalis} unless specifically requested. The CDC does not have universal population-based screening recommendations for trichomoniasis but suggests screening for women at high risk and in high-prevalence settings.

\section*{Why This Test Is Ordered}
\begin{itemize}
  \item Diagnosis of symptomatic vaginal trichomoniasis in women presenting with diffuse, malodorous, yellow-green, frothy discharge, vaginal erythema, dysuria, or pruritus
  \item Screening for \textit{T. vaginalis} in women at high risk: those with multiple sexual partners, history of STIs, commercial sex work, injection drug use, or those residing in high-prevalence settings (correctional facilities, STI clinics)
  \item Evaluation of women with persistent or recurrent vaginal discharge despite treatment for other causes (bacterial vaginosis, candidiasis)
  \item Evaluation of men with persistent non-gonococcal urethritis (NGU) or recurrent urethritis despite treatment for chlamydia, gonorrhea, and \textit{Mycoplasma genitalium}
  \item Screening for \textit{T. vaginalis} in HIV-infected women at entry to care and annually thereafter, as trichomoniasis increases the risk of HIV transmission by increasing genital HIV viral shedding
  \item Evaluation of women with high-risk HPV persistence or cervical dysplasia, given the epidemiologic association between \textit{T. vaginalis} and cervical neoplasia, although screening for \textit{T. vaginalis} for this indication alone is not standard
\end{itemize}

\section*{Primary vs Secondary Test}
NAAT is the primary diagnostic test for trichomoniasis. The wet mount is a primary screening test in settings where NAAT is unavailable or unaffordable, but its low sensitivity means a negative wet mount should be followed by NAAT if clinical suspicion persists. OSOM rapid antigen test is a primary point-of-care alternative. Culture is a secondary test.

\section*{How This Test Is Performed}
Vaginal swab (self-collected or clinician-collected): A Dacron or flocked swab is inserted 2--3 inches into the vagina, rotated for 10--30 seconds, and transferred to the manufacturer-specific transport medium (Aptima vaginal swab transport tube, GeneXpert transport reagent, or standard Amies transport medium if culture is additionally requested). If wet-mount microscopy is to be performed, a separate swab or a drop of vaginal secretion should be placed on a glass slide and examined within 10 minutes (trichomonads lose motility rapidly). For women, vaginal swabs have the highest sensitivity for NAAT (higher than urine). For men, first-catch urine (first 10--20 mL of the urine stream, patient should not have voided for at least 1 hour) or a urethral swab is collected. Urine specimens should be collected in sterile, preservative-free containers and transported to the laboratory within 24 hours at room temperature or within 4 days if refrigerated.

\section*{What Preparation Is Needed}
Patients should not douche, use vaginal creams, suppositories, or spermicides for 24--48 hours before specimen collection. Men should not urinate for at least 1 hour before specimen collection. Avoid specimen collection during menstruation if possible, as blood may interfere with some tests.

\section*{Contraindications and Precautions}
The wet mount's sensitivity is critically dependent on prompt examination (within 10 minutes) by an experienced microscopist. Delays or inexperience dramatically reduce sensitivity. False-negative wet mounts are common (sensitivity 40--60\%); therefore, a negative wet mount should NOT be considered sufficient to exclude trichomoniasis in symptomatic or high-risk patients. NAAT should be performed. Pap smear identification of trichomonads is unreliable and should be confirmed by NAAT. Metronidazole or tinidazole therapy prior to testing will cause false-negative results (all test modalities). Recent douching, intravaginal products, or lubricants may transiently suppress detection. \textit{T. vaginalis} infection is frequently co-existent with bacterial vaginosis (60--80\% co-infection) and other STIs; testing for gonorrhea, chlamydia, HIV, and syphilis should also be performed.

\section*{Risks}
Vaginal or urethral specimen collection is minimally invasive and carries negligible risk. There is no risk from the test itself. The primary clinical risk of undiagnosed trichomoniasis is the enhanced risk of HIV acquisition (2- to 3-fold), adverse pregnancy outcomes (preterm birth, PROM, low birth weight) in pregnant women, and chronic genital inflammation.

\section*{What Happens After the Test}
NAAT results available in 1--4 hours. Wet-mount examination results immediately (but must be performed within 10--20 minutes). OSOM rapid test results in 10 minutes. Confirmed trichomoniasis should be treated with metronidazole 500 mg PO BID for 7 days in women (preferred regimen), or a single 2 g dose (alternative due to lower cure rate ~85\% vs. 95\%). In men, either metronidazole 2 g single dose or 500 mg BID for 7 days is acceptable. Tinidazole 2 g single dose is an alternative but is more expensive. Sexual partners from the preceding 60 days should be treated presumptively (EPT). Patients should abstain from sexual activity until treatment is completed and symptoms resolve. Alcohol consumption should be avoided during metronidazole therapy and for 24 hours after completion (disulfiram-like reaction). Retesting with NAAT at 3 months is recommended for all sexually active women due to high reinfection rates (up to 17\% within 3 months).

\section*{Understanding Results}

\subsection*{Below Normal}
\begin{itemize}
  \item \textbf{NAAT negative (not detected):} No \textit{T. vaginalis} nucleic acid detected. High NPV (>99\%)---essentially excludes infection. However, recent douching, intravaginal products, or metronidazole therapy may reduce organism burden below the detection threshold
  \item \textbf{Wet mount negative (no motile trichomonads seen):} Does NOT exclude trichomoniasis. 40--60\% of infections are missed by wet mount. If clinical suspicion is present (symptoms, known exposure), NAAT should be performed
  \item \textbf{Culture negative:} Does not entirely exclude trichomoniasis, though sensitivity of InPouch culture is 80--90\%. A negative culture should be confirmed by NAAT if suspicion persists
\end{itemize}

\subsection*{Above Normal}
\begin{itemize}
  \item \textbf{NAAT positive for \textit{T. vaginalis}:} Diagnostic of trichomoniasis. Patient and partner should be treated. Asymptomatic infection is common; all NAAT-positive patients warrant treatment regardless of symptoms to prevent sequelae and reduce transmission
  \item \textbf{Wet mount positive (motile, flagellated, pear-shaped trichomonads with jerky motility):} Confirms trichomoniasis. In addition to the presence of trichomonads, the wet mount typically shows an elevated vaginal pH (5--6), increased polymorphonuclear leukocytes (PMNs), and a positive "whiff test" (fishy odor when KOH is added to vaginal discharge, found more commonly in bacterial vaginosis but may be present in trichomoniasis due to co-infection)
  \item \textbf{OSOM rapid test positive:} Confirms trichomoniasis at the point of care, allowing immediate treatment. Sensitivity is 82--90\% compared to NAAT
  \item \textbf{Culture (InPouch) positive:} Confirms viable \textit{T. vaginalis} organisms. The InPouch system is convenient for transport and culture but is being supplanted by NAAT
\end{itemize}

\section*{Normal Results}
\textbf{\textit{Trichomonas vaginalis} not detected} by NAAT. Wet mount: \textbf{No motile trichomonads seen}. \textbf{No \textit{T. vaginalis} isolated} on culture.

\section*{Follow-up Tests}
\begin{itemize}
  \item \textbf{Re-testing with NAAT at 3 months:} Recommended for all sexually active women treated for trichomoniasis due to high reinfection rates (17\% within 3 months)
  \item \textbf{Test of cure (NAAT):} Not routinely recommended for nonpregnant patients treated with the standard regimen. However, if symptoms persist, evaluate for treatment failure (uncommon with metronidazole but resistance is increasing; susceptibility testing available at CDC upon request), reinfection, or co-infection
  \item \textbf{Gonorrhea/Chlamydia NAAT:} Co-testing for gonorrhea and chlamydia is indicated given shared risk factors and high co-infection rates
  \item \textbf{HIV testing and syphilis serology:} For all patients diagnosed with trichomoniasis per STI evaluation guidelines
  \item \textbf{Wet mount with saline and KOH (including clue cells, yeast):} To evaluate for co-existing bacterial vaginosis (clue cells, positive whiff test) and vulvovaginal candidiasis (budding yeast, pseudohyphae), which commonly co-exist
\end{itemize}

\section*{References}
\begin{enumerate}
  \item Workowski KA, Bachmann LH, Chan PA, et al. Sexually transmitted infections treatment guidelines, 2021. MMWR Recomm Rep. 2021;70(4):1-187.
  \item Kissinger P, Mena L, Levison J, et al. A randomized treatment trial: single versus 7-day-dose metronidazole for the treatment of \textit{Trichomonas vaginalis} among HIV-infected women. Lancet Infect Dis. 2018;18(11):1212-1219.
  \item Meites E, Gaydos CA, Hobbs MM, et al. A review of evidence-based care of symptomatic trichomoniasis and asymptomatic \textit{Trichomonas vaginalis} infections. Clin Infect Dis. 2015;61(Suppl 8):S837-S848.
  \item Schwebke JR, Burgess D. Trichomoniasis. Clin Microbiol Rev. 2004;17(4):794-803.
  \item Hobbs MM, Seña AC. Modern diagnosis of \textit{Trichomonas vaginalis} infection. Sex Transm Infect. 2013;89(6):434-438.
\end{enumerate}

\clearpage
\section*{Regulatory Status and Authorization Date}
Regulatory status applies to the specific assay, instrument, manufacturer, and intended use rather than to the test name alone. No single approval date can be assigned to this test category. For a commercial assay, verify the current product in the relevant regulator's database: the U.S. Food and Drug Administration (FDA) clearance, approval, or authorization; India's Central Drugs Standard Control Organisation (CDSCO) licence or permission; the European Union In Vitro Diagnostic Regulation (EU IVDR) CE certificate issued through the applicable conformity-assessment route; or the United Kingdom Medicines and Healthcare products Regulatory Agency (MHRA) registration and UKCA/CE status. Laboratory-developed tests may not have an individual FDA, CDSCO, EU IVDR, or MHRA approval date.
\clearpage
